\subsection{A local linearized solution with a nonlinear obstruction}
\label{sec:oblate_local_obstruction}

The primary--secondary ranks above do not determine whether every retained linearized solution is tangent to a nonlinear solution. A separate local test gives a negative answer for a specified tangent. It uses the original coefficient-Galerkin bulk action, with $N=3$, $\Lambda=0$, and the oblate background at $\epsilon=1/100$ in the radial units below. It does not require a Legendre inverse or identification of this fixed deformation with the sufficiently small neighborhood in Proposition~\ref{prop:oblate_rank}.

\begin{proposition}[Local Lagrangian obstruction]\label{prop:oblate_local_obstruction}
There is an open radial interval $I$ containing $r=1$ and a real stationary retained linearized solution $c_1$ on $\mathbb R_t\times I\times S^2$ about the oblate background such that no family of exact retained bulk solutions on a common spacetime neighborhood has tangent $c_1$ and is $C^2$ in its amplitude in a topology controlling the Euler-equation jets. The solution $c_1$ is the sum of an axisymmetric polar germ and a compactly supported odd coordinate variation.
\end{proposition}

\begin{proof}
First construct the polar germ with $X_0=0$. Write its seven remaining profiles and their radial derivatives as
\begin{align}
 v&=(\lambda_0,\lambda_2,X_2,V_{\mathrm p,2},C_{\mathrm p,2},a_0,a_2)^t,
 & z&=(v,w)^t.
\end{align}
Here $a=\delta h_{tt}$, $w$ is initially an independent seven-component state block, and all profiles are time independent. The shift perturbations vanish. Put
\begin{equation}
 T(r)=\arctan\frac1{10r},\qquad
 \mathcal D=\partial_r-\frac{10}{100r^2+1}\partial_T.
\end{equation}
The archived certificate supplies rational matrices over $\mathbb Q(r,T)$ of sizes $C:6\times14$, $K:14\times14$, and $M:14\times8$. There are no remaining control columns. If $J$ selects the ordered free rows of $M$, exact arithmetic verifies
\begin{equation}\label{eq:oblate_germ_invariance}
 CM=0,\qquad JM=I_8,\qquad (\mathcal DC+CK)M=0.
\end{equation}
The independent profile lift and the original Euler rows also satisfy
\begin{align}\label{eq:oblate_germ_lift}
 (P_vK-P_w)M&=0,\nonumber\\
 (E_0P_v+E_1P_w+E_2P_wK)M&=0.
\end{align}
The matrices $P_v,P_w$ select the upper and lower seven states. Each $E_j$ has all 22 original retained Euler rows after removal of the $X_0$ \emph{column}; the $X_0$ equation is included. The convention is $E_0v+E_1v'+E_2v''$, with derivatives acting on profiles only.

At the physical radius $r_0=1$, rational interval arithmetic using a 25-term alternating arctangent bracket separates all 57 collected proof denominators from zero and certifies a nonzero six-column minor of the saved original $C$. Its zero-based column indices are $(0,1,2,3,4,6)$. These checks imply analytic coefficients and rank six on some smaller open interval $I$; an explicit interval width is not needed here. The original Euler coefficients also have an analytic angular-integral definition for positive $r,\epsilon$. The additional terminal matrices are used as candidate rational functions, whose regularity is checked directly.

Consequently $\operatorname{im}M=\ker C$ on $I$. Set $S=\mathcal DC+CK$. Since $SM=0$, the certified minor gives an analytic matrix $L$ with $S=LC$. Solve the analytic ODE $z'=Kz$ from an initial state in $\operatorname{im}M(r_0)$. Then $h=Cz$ obeys $h'=Lh$ and $h(r_0)=0$, so uniqueness gives $Cz=0$ throughout the germ. Equation~\eqref{eq:oblate_germ_lift} now implies $v'=w$ and all 22 retained Euler equations. No differentiation of $M$ or reduced-coordinate dynamics is required for this forward construction.

Let $e$ denote this stationary polar perturbation and define the discarded response
\begin{equation}\label{eq:oblate_discarded_response}
 D[e]=6\mathscr E^{(1)}_{X,4}[e]-\mathscr E^{(1)}_{C_{\mathrm p},4}[e].
\end{equation}
The coefficient Euler densities $\mathscr E_A$ include the angular average $dx/2$ and the single volume factor $\Sigma$. An independently lifted response row $d$ gives $D[e](r)=d(r)z(r)$. At the same certified radius and arctangent bracket, the numerator of the second component of $dM$ is strictly separated from zero. Taking the second standard unit vector $\mathbf e_2\in\mathbb R^8$ and
\begin{equation}
 z(1)=M\bigl(1,\arctan(1/10)\bigr)\mathbf e_2
\end{equation}
therefore gives $D[e](1)\ne0$. This germ cannot be a pure continuum coordinate variation, whose unrestricted linearized Einstein residual would vanish on the flat oblate spacetime.

It remains to relate this responding germ to nonlinear compatibility. For $f\in C_c^\infty(I)$, the continuum coordinate variation generated by $f(r)x\partial_\phi$ has retained odd coefficients
\begin{equation}
 g_f:\quad C_{\mathrm a,2}=-\frac{2\rho f}{3},\qquad
 V_{\mathrm a,2}=-\frac{f'}3,\qquad W=0.
\end{equation}
It has zero unrestricted linearized Einstein residual. Thus $c_1=e+g_f$ solves all retained linearized equations.

Let $P_c$ be the fixed harmonic projector on the retained \emph{coefficient} fields and $Q_c=1-P_c$. For the angular descriptor $\xi=w_*(r)x\partial_\phi$, with $w_*\in C_c^\infty(I)$, let $F_\xi$ be its unrestricted coefficient transformation. The descriptor satisfies $Q_cF_\xi(c_0)=0$ at the background. The angular Noether identity, after using the retained equations along an exact retained solution, reads
\begin{equation}
 \partial_t C_W[\xi]=-E_{\mathrm{full}}(c)[Q_cF_\xi(c)].
\end{equation}
For a putative family $c(\eta)=c_0+\eta c_1+\eta^2c_2/2+o(\eta^2)$, the background identities $E_{\mathrm{full}}(c_0)=0$ and $Q_cF_\xi(c_0)=0$ give the necessary condition
\begin{equation}\label{eq:oblate_quadratic_compatibility}
 B_\xi[c_1]:=E_{\mathrm{full}}^{(1)}[c_1]
       [Q_cDF_\xi(c_0)c_1]=0.
\end{equation}
Indeed the right side of the Noether identity is $-\eta^2B_\xi+o(\eta^2)$, whereas $C_W[\xi]=0$ on every retained solution. All terms involving $c_2$, including arbitrary time-dependent second-order corrections, multiply one of the two vanishing background factors.

The independently checked discarded coefficient map gives
\begin{equation}
 Q_cDF_\xi(c_0)g_f
 =\frac{8\rho w_*f}{35}e_{X,4}
  -\frac{4\rho w_*f}{105}e_{C_{\mathrm p},4}.
\end{equation}
Reflection under $\phi\mapsto-\phi$ kills the polar self-pairing, and the unrestricted residual of $g_f$ vanishes. Hence
\begin{equation}\label{eq:oblate_mixed_obstruction}
 B_\xi[e+g_f]=\frac4{105}\int_I dr\,\rho w_*fD[e].
\end{equation}
There is no additional radial measure factor. Shrink the support to a subinterval where $D[e]$ has fixed nonzero sign. Choose real smooth compact $f,w_*$ with $fw_*\ge0$ and positive overlap there. Since $\rho>0$, Eq.~\eqref{eq:oblate_mixed_obstruction} is nonzero, contradicting Eq.~\eqref{eq:oblate_quadratic_compatibility}.
\end{proof}

\begingroup\small\raggedright
The exact input matrices, rational interval endpoints, 20 manufactured controls, source hashes, and completed execution record are preserved in \path{reports/oblate-stationary-state-germ-pass-2026-11-28/}. The analytic-domain audit is in \path{reports/oblate-stationary-input-domain-audit-2026-11-28/}; the descriptor map and normalization are in \path{reports/oblate-null-map-check-2026-11-27/} and \path{reports/oblate-stationary-seed-preparation-2026-11-27/}.\par
\endgroup

This result concerns the specified spacetime tangent $e+g_f$. It neither excludes the exact oblate background nor proves an obstruction for $e$ alone or for every representative obtained by adding linearized gauge directions. It supplies no global radial boundary realization, $N=4$ counterpart, final nonlinear Dirac classification, or physical degree-of-freedom count. The eight free initial coordinates above belong to a certified forward stationary family, not a physical mode count.
